\chapter{補助ケース：発明主題レベルの異質解決策型分析}
\chapoverview{本付録では，卒論本文の主分析であるU-Netケースとは別に，発明主題レベルの補助ケースを整理する．具体的には，花王の洗浄・除去系技術から半導体洗浄分野への接続可能性を，特許文書上でどのように確認できるかを示す．この補助ケースは，卒論本文の主張を広げるためではなく，課題・解決手段分離型枠組みが単一手法名に限定されない可能性を示すための発展的材料である．}

\section{位置づけ}
U-Netケースは，明示的な技術手法名を対象とする近接解決策型である．一方，花王の洗浄・除去系技術から半導体洗浄への展開は，洗浄剤組成物，基板洗浄方法，残渣除去技術などの発明主題そのものを対象とする．このため，後者は発明主題レベルの異質解決策型として位置づけられる．

本付録での分析は，技術導入可能性や事業化可能性を直接証明するものではない．あくまで，花王の洗浄・除去系技術資産が，半導体洗浄・基板洗浄・電子部品洗浄領域と特許文書上で接続している可能性を予備的に確認するものである．

\section{データセット}
補助ケースでは，\tabref{tab:kao_dataset}に示す4つのデータセットを用いる．

\begin{table}[htbp]
  \centering
  \caption{発明主題レベル補助ケースのデータセット}
  \label{tab:kao_dataset}
  \begin{tabularx}{\textwidth}{p{24mm}X r r X}
    \toprule
    Dataset & 内容 & Publication数 & Family数 & 役割 \\
    \midrule
    C0-core & 洗浄・除去・リンス系全体 & 61,723 & -- & 参照分野Aの母集団 \\
    C1 & 花王 $\times$ 洗浄・除去・リンス & 398 & 251 & 花王の元技術資産 \\
    S0 & 半導体洗浄分野全体 & 12,183 & 7,073 & 対象分野Bの母集団 \\
    S1 & 花王 $\times$ 半導体洗浄 & 14 & 11 & 花王の半導体洗浄展開候補 \\
    \bottomrule
  \end{tabularx}
\end{table}

\section{予備確認結果}
S1として，花王 $\times$ 半導体洗浄に該当する特許が14 publication / 11 family確認された．S1は件数が小さいため，全件を代表特許確認の対象とできる．S1の目視確認では，6件が半導体基板・シリコンウエハー・基板洗浄として使いやすい代表候補，7件が電子部品・実装・精密基板など周辺領域の洗浄・残渣除去候補，1件がCMP関連の補助的候補として整理された．

この結果は，花王の洗浄・除去系技術資産が，半導体洗浄・基板洗浄・電子部品洗浄領域に特許文書上で接続している可能性を示す．ただし，これは導入可能性や量産適用性を示すものではない．実際の適用可能性は，材料適合性，工程条件，洗浄対象，安全性，コスト，権利範囲などを専門家が確認する必要がある．

\section{卒論本文での扱い}
本補助ケースは，卒論本文では主分析として扱わない．卒論本文の中心は，U-Netを用いた技術手法レベルの近接解決策型分析である．花王ケースは，課題・解決手段分離型枠組みが，単一手法名だけでなく，材料，組成物，処理方法，製造方法などの発明主題レベルにも拡張しうることを示す補助的材料として位置づける．
